% Bagian: second-order-logic
% Bab: syntax-and-semantics
% Subbagian: terms-formulas

\documentclass[../../../include/open-logic-section]{subfiles}

\begin{document}

\olfileid[id]{sol}{syn}{frm}

\olsection{Suku dan \printtoken{P}{formula}}

Seperti dalam logika orde pertama, ungkapan-ungkapan logika orde kedua
dibangun dari kosakata dasar yang memuat \emph{!!{variable}s},
\emph{!!{constant}s}, \emph{!!{predicate}s}, dan kadang-kadang
\emph{!!{function}s}. Dari unsur-unsur tersebut, bersama penghubung logis,
kuantor, dan simbol tanda baca seperti tanda kurung dan koma, dibentuk
\emph{suku} dan \emph{!!{formula}s}. Perbedaannya ialah bahwa selain variabel
untuk objek, logika orde kedua juga memuat variabel untuk relasi dan fungsi,
serta memperbolehkan kuantifikasi atasnya. Jadi, simbol-simbol logis logika
orde kedua adalah simbol-simbol logika orde pertama, dengan tambahan:

\begin{enumerate}
\item Suatu himpunan !!{variable}s relasi orde kedua yang !!{denumerable}s untuk
  setiap aritas~$n$: $\Obj V_0^n$, $\Obj V_1^n$, $\Obj V_2^n$, \dots
\item Suatu himpunan !!{variable}s fungsi orde kedua yang !!{denumerable}s:
  $\Obj u_0^n$, $\Obj u_1^n$, $\Obj u_2^n$, \dots
\end{enumerate}

Sebagaimana kita menggunakan $x$, $y$, dan~$z$ sebagai metavariabel bagi
variabel orde pertama $\Obj v_i$, kita akan menggunakan $X$, $Y$, $Z$, dan
seterusnya sebagai metavariabel bagi~$\Obj V_i^n$, serta $u$, $v$, dan
seterusnya sebagai metavariabel bagi~$\Obj u_i^n$.

\begin{explain}
Simbol-simbol nonlogis suatu bahasa orde kedua ditentukan dengan cara yang
sama seperti bagi suatu bahasa orde pertama: dengan mencantumkan
!!{constant}s, !!{function}s, dan !!{predicate}s dalam bahasa tersebut.

Dalam logika orde pertama, !!{identity}~$\eq$ biasanya disertakan. Dalam
logika orde pertama, simbol-simbol nonlogis suatu bahasa~$\Lang{L}$ sangat
penting agar kita dapat mengungkapkan sesuatu yang menarik. Tentu saja ada
!!{sentence}s yang tidak menggunakan simbol nonlogis, tetapi hanya
dengan~$\eq$ sulit untuk menyatakan sesuatu yang menarik. Dalam logika orde
kedua, karena kita mempunyai persediaan variabel relasi dan fungsi yang tidak
terbatas, kita dapat menyatakan apa pun yang dapat kita nyatakan dalam suatu
bahasa orde pertama, bahkan tanpa persediaan khusus simbol nonlogis.
\end{explain}

\begin{defn}[Suku Orde Kedua]
Himpunan \emph{suku orde kedua} dari~$\Lang L$, $\TrmSOL[L]$, didefinisikan
dengan menambahkan klausa berikut pada \olref[fol][syn][frm]{defn:terms}:
\begin{enumerate}
\item Jika $u$ merupakan variabel fungsi beraritas~$n$ dan $t_1$, \dots,
  $t_n$ merupakan suku, maka $\Atom{u}{t_1, \ldots, t_n}$ merupakan suku.
\end{enumerate}
\end{defn}

\begin{explain}
Jadi, suatu suku orde kedua tampak persis seperti suatu suku orde pertama,
kecuali bahwa pada tempat suku orde pertama memuat !!a{function}~$\Obj{f^n_i}$,
suatu suku orde kedua dapat memuat variabel fungsi~$\Obj{u^n_i}$ sebagai
gantinya.
\end{explain}

\begin{defn}[\usetoken{S}{formula} Orde Kedua]
Himpunan \emph{!!{formula}s orde kedua}~$\FrmSOL[L]$ dari bahasa~$\Lang L$
didefinisikan dengan menambahkan klausa-klausa berikut pada
\olref[fol][syn][frm]{defn:formulas}:
\begin{enumerate}
\item Jika $X$ merupakan variabel predikat beraritas~$n$ dan $t_1$, \dots,
  $t_n$ merupakan suku orde kedua dari~$\Lang L$, maka
  $\Atom{X}{t_1,\ldots, t_n}$ merupakan !!{formula} atomik.

\tagitem{prvAll}{Jika $!A$ merupakan !!a{formula} dan $u$ merupakan variabel
  fungsi, maka $\lforall[u][!A]$ merupakan !!a{formula}.}{}

\tagitem{prvAll}{Jika $!A$ merupakan !!a{formula} dan $X$ merupakan variabel
  predikat, maka $\lforall[X][!A]$ merupakan !!a{formula}.}{}

\tagitem{prvEx}{Jika $!A$ merupakan !!a{formula} dan $u$ merupakan variabel
  fungsi, maka $\lexists[u][!A]$ merupakan !!a{formula}.}{}

\tagitem{prvEx}{Jika $!A$ merupakan !!a{formula} dan $X$ merupakan variabel
  predikat, maka $\lexists[X][!A]$ merupakan !!a{formula}.}{}
\end{enumerate}
\end{defn}

\end{document}
